\textbf{Simulation Free Schrodinger Bridge with reference refinement using normalizing flows.}
Our previous work on Schrodinger bridge or fitting using MMD losses requires simulation of SDEs, which are often computationally intense. In this follow up work, inspired by \cite{bartosh2025sde}, instead of explicitly parameterizing the SDEs by drift and diffusion, we can implicitly parameterize them via time dependent marginal distributions with a diffusion term. Since data only provides information about marginal distributions, matching the model with data should not require accessing the predicted trajectory of the models. With these two observations, we propose that we can parameterize time dependent marginals using a conditional normalizing flow model, conditioned on time and differentiable. With known diffusion structure, this normalizing flow model implicitly defines a flexible class of SDE. And through automatic differentiation with respect to time, we can back out the drift of this implicitly defined SDE. To solve the identifiability problem, we propose to use the KL divergence between this flexible SDE and a reference family as a penalization term in learning objectives. The KL divergence can be written as an expectation of an L2 difference of drifts over marginals at a random time. Since we can sample marginals efficiently without simulating SDEs using the normalizing flow, this KL term can be calculated efficiently. We can perform the same refinement steps with both reference drift and diffusion in the refinement step, resembling \cite{guan2024identifying}. 



\textbf{Assimilating and adaptively collecting data with multiple perturbed dynamics.}
We have shown that domain knowledge can guide inference in single dynamical systems \cite{shen2025multi, berlinghieri2025oh}. Yet unperturbed dynamics cannot fully reveal gene regulation, as many gene expression combinations never occur.  In this follow up work I will focus on inferring shared structure from perturbed stochastic dynamics in Perturb-seq, a technique that takes measurements of many different perturbed systems. Our framework developed in \cite{berlinghieri2025oh} naturally extends single-system inference to perturbation data, treating perturbations as covariates in hierarchical models or assimilating multiple dynamics via generalized Bayes with MMD losses. Beyond inference, it enables uncertainty quantification, supporting adaptive experimental design to identify perturbations that most reduce uncertainty across modalities. 

%I have shown that domain knowledge can guide inference in single dynamical systems \cite{shen2025multi, berlinghieri2025oh}. Yet in biology, unperturbed dynamics cannot fully reveal gene regulation, as many gene expression combinations never occur. Perturb-seq combines genetic perturbations with single-cell RNA sequencing, producing multiple related dynamics that my past work cannot capture. I aim to extend my methods to such multimodal settings: inferring shared structure, analyzing Perturb-seq efficiently, and designing adaptive strategies to select the most informative perturbations. Similar challenges arise in biodiversity surveys, plasma physics, and astronomy, where multiple indirect measurements must be combined. I will study tradeoffs between precision in one modality and generalization across many, and develop principled methods to assimilate and design across them.

%In the short term, I will focus on inferring shared structure from perturbed stochastic dynamics in Perturb-seq. I have connected Schrödinger bridges with Bayesian inference: the bridge emerges as the large-sample limit of a variational Bayes procedure with SDE priors, where volatility estimation acts as an empirical Bayes step. This framework naturally extends single-system inference to multimodal perturbation data, treating perturbations as covariates in hierarchical models or assimilating multiple dynamics via generalized Bayes with MMD losses. Beyond inference, it enables uncertainty quantification, supporting adaptive experimental design to identify perturbations that most reduce uncertainty across modalities. 

